Here $h$ is restricted to the prepared positive odd-sign unit
filter containing one, as in (3.3), so $h\equiv1\pmod8$.
The multiset of dummy labels and its size are fixed before $h$;
its realizing primes are chosen after $h$ and avoid its support.

For the rough companion the curve itself varies:
write $a$ for the bounded dummy product and take $A=E^{ha}$.
Choose $q_0$ so $k_0=aq_0$ is negative, $1$ modulo eight and
a square at every prime of $2Nh$, with the required exclusions.
Then $k_0$ is odd fundamental, prime to $N_{E^h}$, and
\[
          w(E^{hk_0})=w(E^h)\chi_{k_0}(-N_{E^h})=+1.
\]
This proves the sign directly for the actual application.
Apply Hoffstein--Luo to $E^{hk_0}$ with every conductor,
old, $h$, dummy and forbidden prime in its splitting set.
The resulting $k'=k_0d$ adds at most five new primes.
The constants and sizes in the analytic theorem may depend on
this varying curve and the whole finite set; the bound five does not.

\subsection{The even upper transfer}
\begin{proposition}[A uniformly bounded partitionable companion]\label{ex:app-rough-companion}
Fix $E$ and the finitely many prepared even sources, local types and
unit seeds used for (3.3). There are constants $A,C_E$, fixed before
$h$ and before any additional finite prime exclusion set, with the
following property. Let $h$ be positive squarefree in the prepared
odd-sign unit filter containing $1$, prime to its fixed support; in
particular $h\equiv1\pmod8$. There is a negative odd fundamental
discriminant $k'$, prime to $2Nh$ and the additional exclusions,
square at every place dividing $2Nh$, for which $hk'$ is partitionable
in the sense of Definition~\ref{ex:sec7-seed-partition},
$L_*(hk')\ne0$, and
\[
 \omega(k')\le A,\qquad
 v_2L_*(hk')\le2\omega(h)+2\omega(k')+s(hk')+C_E.
\]
Here $A$ is the fixed dummy-label count plus five. Dummy types and
their number are fixed first, but their distinct prime realizations
are selected after $h$ and the exclusions. The proof uses the
one-symbol stationary-address conclusion of
Lemma~\ref{ex:sec7-one-symbol-address}, whose proof is given in
Proposition~\ref{ex:app-simultaneous-address}.
\end{proposition}
\begin{proof}\leavevmode
For one fixed even source write (E.11) as
\[
             v_2L_*(\sigma D)=2v_2a(D)+\kappa_e.
\]
The lower bound (E.9) and a nonzero reference coefficient imply
that the nonempty set of integers
$v_2a(D)-\omega(D)$ has a finite attained minimum $\delta_e$.
A unit of the normalized symbol means precisely equality to this
minimum. It therefore gives
\[
 L_*(\sigma D)\ne0,\qquad
 v_2L_*(\sigma D)=2\omega(D)+C_{e,\rm unit},
 \quad C_{e,\rm unit}=2\delta_e+\kappa_e.                         \tag{F.1}
\]
This constant belongs to the fixed source, not to a curve later
introduced solely to obtain nonvanishing.

For the negative types used in the rough companion, the unsigned
index is $7$ modulo eight. Its total $(-1/D)$ label is nonzero, so
every unit seed is nonempty. A seed of least positive size has no
dependence on internal edge bits: a nonzero edge derivative
would be a smaller nonempty unit. Its labels therefore suffice.
If a list with the same total label contains the first $r-1$
seed anchors and a nonempty remainder, contract that remainder
into the last anchor. Its total label is forced to be the
remaining seed label. This proves partitionability regardless
of how many old $h$-primes are in the remainder.

